\documentclass[12pt]{article}
\usepackage[utf8]{inputenc}
\usepackage[T1]{fontenc}
\usepackage{hyperref}
\title{Urdu Language Resources for K-12 Curriculum Development}
\author{Compiled by Claude Sonnet 5 (claude-sonnet-5) for Andr\'as Kornai}
\date{}
\begin{document}
\maketitle
\section{Introduction}
Urdu (ISO 639-3 urd; Glottocode urdu1245) is an Indo-European language with about 94 million speakers, an official national language of Pakistan and used in education. This search was limited; it did not locate official Pakistani or Indian board textbooks or curriculum documents, which should be sought directly from the relevant boards.
\section{Primary Educational Resources}
\begin{itemize}
\item Urdu Primer, CIIL Mysuru and NCERT (2025 first edition), basal reader covering the alphabet and numerals: \url{https://ciil.gov.in/primers/Urdu_Primer.pdf}
\end{itemize}
Search summaries also mention the heritage series ``Let's Read Urdu'' (readers, workbooks, teacher's manual) but no reliable URL was seen. Generic ``Class 1 syllabus'' and ``Oxford Urdu class 6'' pages appeared but are low-quality SEO pages and are not cited.
\section{Linguistic and Reference Materials}
\begin{itemize}
\item Urdu Lughat (22-volume Urdu-Urdu dictionary, Urdu Dictionary Board): \url{https://en.wikipedia.org/wiki/Urdu_Lughat}
\item Rekhta Dictionary (Urdu, Hindi, English): \url{https://www.rekhtadictionary.com/about-us}
\item Ohio State University Urdu 101 syllabus (university level, not K-12): \url{https://ascnet.osu.edu/storage/request_documents/45/Urdu 101 Syllabus.pdf}
\end{itemize}
\section{Religious and Community Resources}
None found in this search.
\section{Digital and Audio Resources}
\begin{itemize}
\item Rekhta open-access Urdu literature collection, described by McGill Library: \url{https://blogs.library.mcgill.ca/islamicstudieslibrary/?p=2121}
\end{itemize}
\section{Curriculum Development Context}
Urdu is a language of instruction, so the main task is gathering official board textbooks and curricula (Pakistan national and provincial boards, Indian state boards and NCERT), which this search did not reach. The CIIL/NCERT primer suits early literacy; the dictionaries and Rekhta literature support upper grades.
\end{document}